\documentclass[12pt]{article}
\usepackage[utf8]{inputenc}
\usepackage{amsmath,amssymb,amsthm}
\usepackage{algorithm}
\usepackage{algpseudocode}
\usepackage{geometry}
\geometry{a4paper, margin=1in}

\newtheorem{theorem}{Theorem}
\newtheorem{lemma}[theorem]{Lemma}
\newtheorem{corollary}[theorem]{Corollary}
\newtheorem{proposition}[theorem]{Proposition}
\theoremstyle{definition}
\newtheorem{definition}[theorem]{Definition}

\title{Project Euler Problem 808: Reversible Prime Squares}
\author{}
\date{}

\begin{document}
\maketitle

\section{Problem Statement}

Both 169 and 961 are the square of a prime. 169 is the reverse of 961. A number is called a \textbf{reversible prime square} if:
\begin{enumerate}
    \item It is not a palindrome.
    \item It is the square of a prime.
    \item Its digit-reversal is also the square of a prime.
\end{enumerate}

Find the sum of the first 50 reversible prime squares.

\section{Mathematical Analysis}

\subsection{Pairing Property}

\begin{theorem}
If $n = p^2$ is a reversible prime square, then $\mathrm{rev}(n) = q^2$ is also a reversible prime square.
\end{theorem}

\begin{proof}
Let $n = p^2$ be a reversible prime square. Then $n$ is not a palindrome, $p$ is prime, and $\mathrm{rev}(n) = q^2$ for some prime $q$. We need to show $q^2$ is also a reversible prime square:
\begin{itemize}
    \item $q^2$ is the square of a prime $q$. \checkmark
    \item $\mathrm{rev}(q^2) = \mathrm{rev}(\mathrm{rev}(p^2)) = p^2$ (assuming no leading zeros), which is the square of the prime $p$. \checkmark
    \item Since $n = p^2$ is not a palindrome, $p^2 \neq q^2$, so $q^2$ is not a palindrome either. \checkmark
\end{itemize}
\end{proof}

\subsection{Search Space}

For primes $p$ up to $10^8$, the squares $p^2$ can be up to $10^{16}$. This provides a sufficient range to find at least 50 reversible prime squares.

\subsection{Digit Constraints}

For $p^2$ to satisfy the conditions:
\begin{itemize}
    \item The leading digit of $p^2$ must be a valid last digit of a perfect square, i.e., one of $\{0, 1, 4, 5, 6, 9\}$.
    \item The last digit of $p^2$ cannot be 0 (to avoid leading zeros in the reversal).
\end{itemize}

\section{Algorithm}

\begin{algorithmic}[1]
\State Sieve all primes up to $N = 10^8$
\For{each prime $p$ in the sieve}
    \State $s \gets p^2$
    \If{$s$ is a palindrome} \textbf{continue} \EndIf
    \State $r \gets \mathrm{rev}(s)$
    \State $q \gets \lfloor\sqrt{r}\rfloor$
    \If{$q^2 \neq r$} \textbf{continue} \EndIf
    \If{$q$ is not prime} \textbf{continue} \EndIf
    \State Add $s$ to results
\EndFor
\State Sort results and sum the first 50
\end{algorithmic}

\section{Complexity Analysis}

\begin{itemize}
    \item \textbf{Time}: $O(N \log \log N)$ for the sieve, $O(\pi(N))$ for the main loop.
    \item \textbf{Space}: $O(N)$ for the prime sieve.
\end{itemize}

\section{Answer}

\[
\boxed{3807504276997394}
\]

\end{document}
